\documentclass[12pt]{article}
\usepackage[utf8]{inputenc}
\usepackage[T1]{fontenc}
\usepackage{hyperref}
\title{Suundi Language Resources for K-12 Curriculum Development}
\author{Compiled by Claude Sonnet 5.5 (claude-sonnet-5-5) for Andr\'as Kornai}
\date{}
\begin{document}
\maketitle

\section{Introduction}
Suundi (Kisuundi, Kissoundi, Sundi; ISO 639-3 sdj; Glottocode suun1239) is a Bantu language of the Kongo group, spoken mainly in the Republic of the Congo (and, per the reference description, Angola). Wikipedia reports about 120,000 speakers (2000). Some Congolese sources treat Kisundi as a dialect of Kikongo. Reported dialects include Kisuundi kya Mayama and Kisuundi kya Ngwaari. Very little published material was found.

\section{Primary Educational Resources}
No primers, textbooks or school programmes were found.

\section{Linguistic and Reference Materials}
\begin{itemize}
\item Wikipedia article on the Suundi language: \url{https://en.wikipedia.org/wiki/Suundi_language}
\item Wikipedia article on Sundi: \url{https://en.wikipedia.org/wiki/Sundi}
\item Joshua Project language page for sdj, with alternate names and dialects: \url{https://joshuaproject.net/index.php/languages/sdj}
\item Glottolog bibliographic record for F.\ Lumwamu (1977), \emph{Lexique kikongo-fran\c{c}ais, 3: kiladi-kisundi-kikongo} (Brazzaville, 151 pp.), a possible Kisundi lexicon. Glottolog codes it under Laari, so its coverage of Suundi is uncertain: \url{https://glottolog.org/resource/reference/id/82521}
\end{itemize}

\section{Religious and Community Resources}
\begin{itemize}
\item Joshua Project profile of the Sundi people in the Republic of the Congo. It lists Bible translation as started, with no portions, New Testament or full Bible available: \url{https://joshuaproject.net/index.php/people_groups/15122}
\end{itemize}

\section{Digital and Audio Resources}
No audio, video, corpus or software resources specific to Suundi were found.

\section{Curriculum Development Context}
French is the language of schooling in the Republic of the Congo, with Kikongo-related varieties (including Kituba and Lari) used as lingua francas. No Suundi school materials or standard orthography were found. The language's relation to Kikongo, and whether it is a separate language or a dialect, should be settled with the community before any curriculum work. The Lumwamu lexicon is the most concrete lead.
\end{document}